% Transcription — fonds Alexandre Grothendieck, shelfmark 54, batch 1 (pages 1-7).
% Established from the facsimile archives/batches/54/batch-01.pdf (a mirror of
% https://grothendieck.umontpellier.fr/54.pdf), per the skill
% transcribe-grothendieck. The folder is seven pages and this is its only batch.
% Page 1 is a cover in his hand, naming the run; pages 2, 4 and 6 carry that run
% — the duality between the cohomology of an abelian variety and that of its
% dual, taken first through the universal extension by a vector group and then,
% in characteristic p, through Frobenius and Verschiebung. Pages 3 and 5 are two
% leaves of a different document altogether: a typescript on flatness and
% freeness, numbered 4.1 and following, which he has gone over in ink, striking
% a hypothesis out of one statement and renaming another. Following the practice
% settled on folder 48, only the statements his hand actually works on are set
% here; the body of the typescript is printed text and not a hand, and is not
% re-set. Page 7 is a Bourbaki typescript on chambers and walls — the recto of
% the leaf carrying page 6 — unrelated and untouched by him: it does not appear,
% and the gap in the numbering is the record that it was skipped.
% Pass: Opus 5 (claude-opus-5), 2026-08-30 — first pass, unchecked against the pages by a human.
\documentclass[11pt,a4paper]{article}
% Le préambule commun à toutes les transcriptions du fonds Grothendieck.
%
% Il tient en une idée : **l'incertitude se compose, elle ne se résout pas.**
% Une transcription qui lisse un mot douteux a détruit la seule chose pour
% laquelle on la faisait — un lecteur doit pouvoir savoir, à chaque endroit,
% ce qui était sur le feuillet et ce qui a été ajouté. Les six macros qui
% suivent sont donc l'appareil critique, et non de la décoration.
%
% Les mêmes commandes sont reconnues par `scripts/render.mjs`, qui produit la
% vue de lecture du site. Une macro ajoutée ici doit l'être aussi là-bas,
% faute de quoi le rendu échouera bruyamment — ce qui est voulu : un rendu
% silencieusement faux est pire qu'un rendu absent.

\ProvidesPackage{grothendieck}[2026/08/08 Transcriptions du fonds Grothendieck]

\usepackage{amsmath,amssymb,amsthm}
\usepackage{mathrsfs}
\usepackage{stmaryrd}
% Les diagrammes commutatifs sont partout dans ce fonds : c'est la langue
% même de la géométrie algébrique de Grothendieck, et une transcription qui
% les rendrait en prose ne transcrirait rien.
\usepackage{tikz-cd}
% Français par défaut — c'est la langue des feuillets — mais l'anglais est
% chargé aussi : la traduction et le résumé sont des documents anglais, et la
% typographie française (espaces avant les deux-points) y serait une faute.
% Ces documents basculent par \selectlanguage{english} en tête de corps.
\usepackage[english,french]{babel}
\usepackage{fontspec}
\usepackage{geometry}
\geometry{a4paper,margin=25mm,marginparwidth=28mm}
\usepackage{marginnote}
\usepackage{xcolor}
% ulem, not soul. `soul`'s \st and \ul are built on a character-by-character
% scan that XeLaTeX's font machinery breaks: the document fails with an
% undefined control sequence rather than degrading. ulem does the same two jobs
% and is Unicode-engine safe. `normalem` stops it hijacking \emph, which the
% transcriptions use for Grothendieck's own emphasis.
\usepackage[normalem]{ulem}
\usepackage{hyperref}
\hypersetup{colorlinks=true,linkcolor=black,urlcolor=black,pdfborder={0 0 0}}

% --- Métadonnées du lot -----------------------------------------------------
% Reprises telles quelles par le script de rendu, qui les affiche en tête de
% la vue de lecture. Elles fixent ce que le fichier prétend couvrir : sans
% elles, une transcription flotte sans point d'attache dans un fonds de seize
% mille feuillets.

\newcommand{\folder}[1]{\gdef\grothFolder{#1}}
\newcommand{\batch}[1]{\gdef\grothBatch{#1}}
\newcommand{\leaves}[2]{\gdef\grothFirst{#1}\gdef\grothLast{#2}}
\newcommand{\foldertitle}[1]{\gdef\grothFoldertitle{#1}}
\gdef\grothFolder{?}\gdef\grothBatch{?}\gdef\grothFirst{?}\gdef\grothLast{?}\gdef\grothFoldertitle{}

% --- Le résumé --------------------------------------------------------------
%
% Il ouvre la lecture modernisée, sous le titre « Résumé », et il est composé
% en petit. Ce n'est pas un ornement : c'est le seul passage du document qui
% ne soit pas des mathématiques, et un lecteur doit pouvoir le reconnaître
% avant de l'avoir lu — puis le sauter s'il connaît déjà le sujet, ou s'y
% arrêter s'il ne le connaît pas. Le filet qui le suit marque la reprise du
% corps.

\newenvironment{resume}
  {\small\setlength{\parskip}{0.45em}}
  {\par\medskip\hrule\medskip}

% --- Mots-clés ---------------------------------------------------------------
%
% En anglais, en clôture du résumé de la lecture modernisée : le vocabulaire
% moderne sous lequel chercher ce que les feuillets construisent. Le manifeste
% du site (`npm run manifest`) les extrait pour étiqueter la cote entière —
% cette ligne est la source unique des tags, il n'y a pas de fichier à côté.
\newcommand{\keywords}[1]{\par\smallskip\noindent
  {\footnotesize\textsc{Keywords} --- \textit{#1}}\par}

% --- L'appareil critique ----------------------------------------------------

% Le numéro de feuillet, en marge. C'est l'ancre qui permet de régler un
% désaccord sur une formule en regardant *un* feuillet plutôt qu'une chemise
% de six cents. Le site s'en sert aussi pour tourner les pages du fac-similé
% quand on fait défiler la transcription.
\newcommand{\leaf}[1]{%
  \marginnote{\footnotesize\color{gray!70}\textsf{#1}}%
  \label{leaf:#1}\ignorespaces}

% Illisible : on ne devine pas. Un mot inventé qui se lit comme les autres est
% le pire résultat possible.
\newcommand{\ill}{\textcolor{red!60!black}{[\,\ldots\,]}}

% Lecture douteuse : le mot est donné, et signalé comme incertain.
\newcommand{\uncertain}[1]{\uline{#1}}

% Ajout de l'éditeur — une lettre manquante, un mot sous-entendu.
\newcommand{\add}[1]{\textcolor{blue!50!black}{[#1]}}

% Biffé par Grothendieck lui-même. Ce qu'il a rayé fait partie du document :
% une rature dit souvent où la pensée a changé de direction.
\newcommand{\struck}[1]{\sout{#1}}

% Note du transcripteur, sur l'état matériel du feuillet.
\newcommand{\note}[1]{\par\smallskip{\small\color{gray!60!black}\textit{[#1]}}\par\smallskip}

% Note marginale de Grothendieck, distincte du corps du texte.
\newcommand{\marginal}[1]{\par\smallskip\noindent\hspace{1em}%
  {\small\color{gray!60!black}#1}\par\smallskip}

% --- Titre ------------------------------------------------------------------

\AtBeginDocument{%
  \begin{center}
    {\large\bfseries Fonds Alexandre Grothendieck --- cote n\textsuperscript{o}~\grothFolder}\\[2pt]
    {\itshape\grothFoldertitle}\\[4pt]
    {\small feuillets \grothFirst--\grothLast\ (lot \grothBatch)}\\[2pt]
    {\footnotesize\color{gray!60!black}Transcription établie d'après le fac-similé numérisé
     par l'Université de Montpellier.\\
     Les lectures incertaines sont soulignées, les illisibles notées [\,\ldots\,],
     les ajouts entre crochets.}
  \end{center}
  \vspace{1em}}

\endinput


\folder{54}
\batch{1}
\pages{1}{7}
\dating{[à partir de 1971]}
\watermark{Édition de démonstration}
\foldertitle{Dualité en cohomologie des V. A [variétés abéliennes] : notes manuscrites (s.d.)}

\begin{document}

\section*{Dualité en Cohomologie de V.A.}

\note{titre de sa main, en haut à droite du feuillet de garde, qui ne porte rien
d'autre ; c'est lui qui nomme la suite et il ne reçoit pas de numéro de page}

\page{2}

$A$ V.A., $B$ duale.

\[
  H^{*}(A) \;=\; \underline{\underline{H}}{}^{*}\!\left(A,\ \underline{\underline{\Omega}}{}^{*}_{A}\right)
\]

\note{le $H$ et le $\Omega$ sont l'un et l'autre soulignés deux fois ; glosé à
droite de la formule, sur la même ligne : « \uncertain{Coh.} de
Hodge--DeRham »}

\[
  t_B = H^1(A, \mathcal{O}_A), \qquad
  t_A = H^0(A, \Omega^1_A)^{\vee} \simeq H^n(A, \Omega^{n-1}_A) .
\]

\emph{Autre interprétation :} Considérons l'\emph{extension universelle} de $A$
par un groupe vectoriel \add{$V$}, \struck{elle} \ill{} \struck{\ill{}} par
$\operatorname{Hom}\!\left(H^1(A,\mathcal{O}_A)^{\vee},\, V\right)$, \ill{}
l'extension universelle

\note{la phrase est reprise deux fois : « $V$ une telle extension est » est
écrit en interligne au-dessus de « un groupe vectoriel », et la formule
$\operatorname{Hom}(H^1(A,\mathcal{O}_A)^{\vee}, V)$ est insérée au-dessus d'un
$V$ biffé, à la suite de $H^1(A,\mathcal{O}_A)$. Les mots de liaison qui
tenaient les deux états ensemble ne se laissent pas lire}

\[
  0 \longrightarrow W(\check{t}_B) \longrightarrow E \longrightarrow A \longrightarrow 0
\]

\note{le caron sur $t_B$ est sa marque du dual ; il ne l'explique pas ici}

et \uncertain{passant} aux espaces tangents on trouve une extension

\[
  0 \longrightarrow \check{t}_B \longrightarrow t_E \longrightarrow t_A \longrightarrow 0 .
\]

Je dis que cette extension n'est autre que $S_B$, (ce qui donne en prenant
$t_E = H^1(B)$) et que de plus elle est \emph{transposée} de $S_A$, ce qui
donne

\begin{quote}
$H^1(A)$ et $H^1(B)$ sont \uncertain{duals} l'un de l'autre
\end{quote}

\note{encadré sur la page. $S_A$ et $S_B$ ne sont pas définis dans ce qui
précède ; il les suppose connus}

Plus fin, il \uncertain{faut} prouver que $H^{*}(A)$ et $H^{*}(B)$

\page{3}

\note{feuillet d'un autre document, que le dossier a reçu comme support : un
tapuscrit sur la platitude et la liberté, numéroté 4.1 et suivants, écrit à la
première personne (« Je m'aperçois à l'instant d'une amélioration dans une autre
direction de 4.1., dont je vais reprendre l'énoncé »). On n'en reprend ici que
les énoncés sur lesquels sa plume travaille}

\textbf{Proposition 4.1.} Soient $A$ un anneau, $I$ un idéal
\struck{nilpotent} de $A$, $(A_i)_{i \in I}$ une famille de $A$-algèbres telle
que l'intersection des noyaux des homomorphismes $A \to A_i$ soit réduite à
$0$, $M$ un $A$-module. On suppose $M \otimes_A A/I$ libre sur $A/I$ (resp.\
plat sur $A/I$, et l'ensemble d'indices $i$ fini), et pour tout $i$,
$M \otimes_A A_i$ libre (resp.\ plat) sur $A_i$. Alors \add{$\forall n$,} $M$
\add{$\otimes_A (A/I^{n+1})$} est \struck{plat sur $A$} libre (resp.\ plat) sur
$A/I$\add{$^{n+1}$}.

\note{c'est là toute l'intervention de sa main sur ce feuillet, et elle est
mathématique : l'hypothèse de nilpotence tombe, et la conclusion, qui portait
sur $A$ lui-même, ne porte plus que sur les quotients $A/I^{n+1}$, pour tout
$n$. Le tapuscrit demandait $I$ nilpotent précisément pour passer de $A/I$ à
$A$ ; la phrase suivante, restée telle quelle, s'appuie encore dessus}

\textbf{Corollaire.} \struck{Supposons} $A$ \uncertain{anneau}
\struck{\uncertain{noethérien}}, \ill{} $\left(\bigcap I_i = 0\right)$, \ill{}
$M$ \ill{} $A$-\uncertain{module} \ill{}

\note{brouillon écrit à la main au bas du feuillet, sous un grand crochet, et
biffé ligne après ligne ; on n'en tire que les mots ci-dessus. On y relit
$M \otimes_A A/I_i$ et « $M$ plat / $A$ », mais l'énoncé ne se laisse pas
reconstituer et n'est pas reconstitué ici}

\page{4}

sont \uncertain{duals} l'un de l'autre, et que $H^p(A, \Omega^q_A) \to$ soit le
dual de $H^q(B, \Omega^p_A)$.

\note{l'indice du second $\Omega$ est écrit $A$ ; la symétrie de l'énoncé
demande $B$. La lecture est nette et n'est pas corrigée}

(Ce dernier fait, \struck{\ill{}} : $H^{*}(A,\Omega^{*}_A) =
\bigwedge\left(H^{*}(A,\Omega^{*}_A)\right)^{(1)}$ est immédiate.)
\struck{De plus,} pour avoir une dualité entre $H^{*}(A)$ et $H^{*}(B)$, il
\struck{suffirait} \ill{} structure commune de \uncertain{les définir} entre
$H^{*}(A)$ et $H^{*}(B)$, ce qui ne saurait offrir de difficultés\ldots

\note{$(1)$ marque la composante de degré $1$ : l'anneau de cohomologie de
Hodge d'une variété abélienne est l'algèbre extérieure sur cette composante.
La phrase qui suit est reprise en cours de route et ses mots de liaison ne se
lisent pas tous}

Je \ill{} bien \ill{} de conjuguer ces extensions et dualités avec ce qui donne
la \emph{dualité de Cartier}. Considérons (le cas \uncertain{strict} : $p > 0$)
${}_pA$ qui est dual au sens Cartier de ${}_pB$. On a (ainsi
$F^A = \operatorname{gr}(t_A) \subset {}_pA$), de degrés respectifs
$p^{\,\struck{2d}\,d}$ et $p^{2d}$, où $d = \dim A$.

\note{l'exposant du premier $p$ est corrigé de $2d$ en $d$ : c'est le rang de
$F^A$, contre $p^{2d}$ pour ${}_pA$ tout entier}

\[
  \left.
  \begin{array}{l}
    F^A \subset {}_pA \\[2pt]
    F^B \subset {}_pB
  \end{array}
  \right\} \quad \text{duals l'un de l'autre}
\]

Je dis que $F^A$ et $F^B$ sont \emph{orthogonaux} l'un de l'autre.
\struck{Noter que $F^B$ \ill{}}

\page{5}

\note{second feuillet du même tapuscrit ; même règle qu'au feuillet 3}

\add{Proposition A} \struck{\textbf{Corollaire 4.1. ter}} Soit $f : X \to Y$ un
morphisme de présentation finie, avec $Y$ \struck{affine} \add{$\neq \emptyset$},
$F$ un Module de présentation finie sur $X$. Alors, quitte à remplacer $Y$ par
un ouvert non vide convenable, on a ce qui suit : il existe un sous-schéma fermé
$Z$ de $Y$, \add{de présentation finie sur $Y$,} tel que pour tout changement de
base $Y' \to Y$, la condition nécessaire et suffisante pour que $F'$ soit plat
par rapport à $Y'$ est que $Y' \to Y$ se factorise en $Y' \to Z \to Y$.

\note{trois interventions de sa main, et la première est la plus lourde de
conséquence : l'énoncé cesse d'être un corollaire de 4.1 pour devenir une
proposition autonome, « Proposition A ». Le feuillet 3 se sert déjà de ce nom
(« En vertu de \emph{Prop.\ A} »), mais là c'est la machine qui l'a frappé,
par-dessus une référence à 4.1 biffée au ruban ; le tapuscrit se citait donc
sous un nom que son propre titre ne portait pas, et c'est la plume qui met le
titre à jour. L'hypothèse « $Y$ affine » tombe au profit de
« $Y \neq \emptyset$ », et $Z$ reçoit d'être de présentation finie sur $Y$}

De plus, (quitte à remplacer $Y$ par un ouvert non vide plus petit), si $Y$ est
affine et $X$ est réunion finie d'ouverts affines $U_i$, alors on a ceci : si la
condition précédente est vérifiée pour $Y'$, et si $Y'$ est d'anneau $A'$
affine, alors pour tout $i$, $\Gamma(U_i, F')$ est un $A'$-module \emph{libre}.

\page{6}

\ill{} sur le diagramme de deux suites transposées l'une de l'autre

\begin{tikzcd}
A \arrow[r, "F"] & A^{(p)} \arrow[r, "V"] & A \\
B & B^{(p)} \arrow[l, "V"] & B \arrow[l, "F"]
\end{tikzcd}

\note{le troisième sommet de la ligne du haut est tracé comme un $V$ nu, sans
barre ; on le lit $A$ par la symétrie avec la ligne du bas, qui est complète, et
la lecture reste douteuse}

On a donc ainsi la suite exacte

\[
  0 \longrightarrow \operatorname{Gr}(t_A) \longrightarrow {}_pA
    \longrightarrow D\!\left(\operatorname{Gr}(t_B)\right) \longrightarrow 0
\]

extension canonique de $D(F^B)$ par $F^A$.

Mais d'ailleurs \ill{} \struck{\ill{}} et correspond \ill{} une extension \ill{}
les \uncertain{correspondances} de ${}_pA$ ??

\note{la fin du feuillet est un brouillon rapide, entrecoupé d'une insertion
entourée d'un trait de plume que l'encre ne rend pas ; les formules portent, les
phrases entre elles non. Les « ?? » sont de lui et ferment la page}

\end{document}
